% Pista ana-25j-q4a · Anàlisi
% Procedència: PAU Catalunya 2025 · Convocatòria de juny · Sèrie 1 · Exercici 4 — Opció A
\documentclass[11pt,a4paper]{article}
\usepackage{pista}
\pistainfo{Catalunya 2025}{Convocatòria de juny}{Sèrie 1}

\begin{document}

\section*{\textcolor{blauPista}{Exercici 4 --- Opció A}}

\noindent La vela major d'un veler té forma semiparabòlica i està delimitada per les gràfiques de $f(x) = -x^{2} + 25$, $y = 0$ i $x = 0$, tal com s'indica a la figura següent:

\noindent\textit{(A la figura: la vela és la regió del primer quadrant compresa entre la paràbola i els eixos de coordenades.)}

\noindent La vela té dues parts separades per la recta $y = 9$. Per a construir-la, s'empra un teixit de niló a la part superior, que costa $50\,\text{\euro}/\text{u}^{2}$, i un teixit de polièster a la part inferior, que costa $70\,\text{\euro}/\text{u}^{2}$. Calculeu el cost total del material que es necessita per a construir aquesta vela. \hfill\textnormal{[2,5~punts]}

\pista{Primer troba on la paràbola talla la recta $y = 9$ i l'eix d'abscisses (de cada equació, queda't amb la solució positiva): seran els límits d'integració. Planteja cada àrea com una integral definida respecte de $x$, i recorda que la regió inferior té una part rectangular (fins al tall amb $y = 9$) i una altra limitada per la corba (des d'aquest punt fins al tall amb l'eix $OX$). Per a la regió superior, calcula l'àrea total sota la paràbola i resta-hi la de la regió inferior. El cost final és la suma de cada àrea multiplicada pel preu corresponent.}

\vspace{0.6em}

\end{document}
